\section{涉外刑事诉讼程序}
\begin{enumerate}
  \item 本身不是法考的重点。
  \item 涉及外国人的 \hz 需要在国外进行刑诉活动的一些特别的方式方法和步骤。
  \item \zd{范围（扩大解释）}。
        \begin{enumerate}
          \item 可能需要到国外去调查取证。
          \item 不仅限于涉外案件的程序，也包括非涉外案件。
        \end{enumerate}
  \item 法院。
        \begin{enumerate}
          \item 主权原则。适用中国刑事法律和信守国际条约相结合。
          \item \zd{基层法院}管辖。
          \item 使用中国通用的语言文字（庭审，文书等）进行诉讼。
                \begin{enumerate}
                  \item 提供找翻译。拒绝翻译的，书面声明记录在案，\zd{必要时应当}录音录像。
                  \item 诉讼文书以中文文本为准，应当附有外文的译本。
                  \item 判决书为中文文本。
                \end{enumerate}
          \item \zd{强制措施要通知}有关国家驻华\zd{使领馆}、同级政府\zd{外事部门、主管部门和当事人}。
          \item 限制出境（扣留护照 \hz 其他出入境证件）。必须到庭的\zd{证人可以要求暂缓出境}。
          \item \zd{没有委托辩护人，可以不提供强制法律援助}。
        \end{enumerate}
  \item 当事人。
        \begin{enumerate}
          \item 国民待遇原则。享有中国法定的诉讼权利并承担诉讼义务。
          \item \zd{国籍判断}。入境时\zd{持用}的有效证件 \yjt 公安 \hz 有关国家驻华使领馆出具的\zd{证明} \yjt 国籍不明。
          \item \bs{探视、会见、旁听}向受案法院所在地\bs{高院}提出。
          \item 可以\zd{拒绝不要翻译}，但需要向法院\zd{书面声明}。
          \item 委托律师只能委托中国律师护代理。
          \item 跨国委托书。同个人提供来自境外的证据材料。
                \begin{enumerate}
                  \item \zd{一个证明}。所在国公证机关的证明。
                  \item \zd{两个认证}。所在国中央外交主管机关 \hz 授权机关的认证；中国驻该国使领馆的认证。
                \end{enumerate}
        \end{enumerate}
\end{enumerate}
